\chapter{Friedmann Dynamics: Matter, Radiation, and Vacuum Energy}

\lectureoverview{The kinematics of an expanding universe can be stated in
one line: physical separations grow with a scale factor \(a(t)\). The problem
now is to determine that function. We shall recover the flat-space Friedmann
equation from Newtonian energy conservation, solve it first for ordinary
matter, and then see why radiation and vacuum energy lead to entirely
different cosmic histories.}

\section{The geometry to be explained}

At a fixed cosmic time, the large-scale universe appears homogeneous and
isotropic. Galaxies, clusters, and superclusters are certainly real features,
but on scales of hundreds of millions or billions of light-years their
distribution becomes remarkably uniform. This empirical regularity is called
the cosmological principle. The name does not make it a theorem; the force of
the principle comes from observation.

There is a second empirical fact. Within present accuracy, large spatial
triangles behave as Euclidean triangles. The comparison with surveying a
football field is useful. If tightly stretched lines make triangles whose
angles add to \(180^\circ\), one has not proved that Earth is globally flat.
One has shown that Earth's curvature radius is large compared with the field.
Cosmic surveying has the same logical form: space may have a very large
curvature radius, but over the observable region it is well approximated as
flat.

The reference geometry is Minkowski spacetime,
\begin{equation}
  d\tau^2=dt^2-\left(dx^2+dy^2+dz^2\right).
  \label{eq:l09-minkowski}
\end{equation}
The bracketed expression is ordinary three-dimensional Pythagoras. In general
relativity, by contrast, the coefficients of the coordinate differentials
are fields: they may depend on position and time, and mixed terms may also
occur. Those metric coefficients are the relativistic successors of the
Newtonian gravitational field.

\begin{figure}[H]
  \centering
  \includegraphics[width=0.72\textwidth]{lecture_09_figure_01.png}
  \caption{The Minkowski proper-time element written at the beginning of the
  lecture. Source frame at 00:01:50.}
  \label{fig:l09-minkowski-board}
\end{figure}

For a spatially flat expanding universe, homogeneity and isotropy leave one
time-dependent scale factor:
\begin{equation}
  d\tau^2=dt^2-a^2(t)\left(dx^2+dy^2+dz^2\right).
  \label{eq:l09-flat-frw}
\end{equation}
We use units in which \(c=1\), so time and length have the same dimensions.
There are two consistent conventions. One may give \(x,y,z\) dimensions of
length and take \(a\) dimensionless, or treat \(x,y,z\) as dimensionless
labels and let \(a\) carry the dimensions of length. The lecture uses the
second convention because it makes the physical meaning of \(a(t)\) visible.

\begin{figure}[H]
  \centering
  \includegraphics[width=0.72\textwidth]{lecture_09_figure_02.png}
  \caption{The flat expanding metric, with the scale factor multiplying the
  Euclidean spatial line element. Source frame at 00:09:12.}
  \label{fig:l09-expanding-metric-board}
\end{figure}

Imagine coordinates painted on a rubber sheet. The coordinate marks remain
attached to the sheet while the physical distance between them grows. Thus a
coordinate difference such as "from \(x=0\) to \(x=1\)" is only a pair of
labels. The distance represented by that interval is \(a(t)\). Finding the
equation that governs this scale factor is the dynamical problem of the
lecture.

\section{Closed, flat, and open spatial geometries}

Flatness is an observational simplification, not a consequence of
homogeneity. There are three homogeneous and isotropic spatial geometries:
positive, zero, and negative curvature. To understand the positive case, it
helps to be precise about dimension.

The locus
\begin{equation}
  X_1^2+X_2^2=R^2
\end{equation}
is a circle, or one-sphere \(S^1\). Its filled interior is a disk, also called
a two-ball. Adding one embedding coordinate gives
\begin{equation}
  X_1^2+X_2^2+X_3^2=R^2,
\end{equation}
the ordinary two-sphere \(S^2\). Adding a fourth embedding coordinate gives
\begin{equation}
  X_1^2+X_2^2+X_3^2+X_4^2=R^2,
  \label{eq:l09-three-sphere}
\end{equation}
a three-sphere \(S^3\), which is itself a three-dimensional space. We cannot
picture its four-dimensional embedding, but the intrinsic geometry does not
need that embedding in order to exist.

\begin{figure}[H]
  \centering
  \includegraphics[width=0.72\textwidth]{lecture_09_figure_03.png}
  \caption{The embedding equation used to introduce the three-sphere. Source
  frame at 00:16:40.}
  \label{fig:l09-three-sphere-board}
\end{figure}

Every point of an ideal sphere is geometrically equivalent to every other,
and every direction at a point is equivalent. A sphere is therefore both
homogeneous and isotropic. If its radius is multiplied by \(a(t)\), the whole
space expands or contracts without selecting a center on the sphere. Tracing
such a universe backward does not place all matter at one special location in
an exterior void. The entire space becomes small, so every pair of points
becomes close. The Big Bang is not an explosion from one point of space.

On a positively curved sphere, geodesic triangles have angle sum greater than
\(180^\circ\). Flat triangles have angle sum exactly \(180^\circ\). In a
homogeneous negatively curved geometry, represented imperfectly by a saddle,
the angle sum is less than \(180^\circ\). These intrinsic tests distinguish
the three cases.

\begin{figure}[H]
  \centering
  \includegraphics[width=0.72\textwidth]{lecture_09_figure_07.png}
  \caption{Positive, zero, and negative curvature sketched through their
  characteristic triangle shapes. Source frame at 00:37:30.}
  \label{fig:l09-curvature-board}
\end{figure}

A perfectly flat local metric does not by itself prove that space is
topologically infinite. Compact flat topologies are mathematically possible,
although the lecture adopts the simplest infinite Euclidean model.
\editorialnote{This global-topology qualification is not needed for any of the
subsequent local dynamics.}

\section{Comoving labels and Hubble's law}

Return now to the expanding flat sheet and draw a coordinate grid on it. In
comoving coordinates, the grid expands with the large-scale matter flow. A
galaxy, or more accurately a sufficiently coarse-grained cluster of galaxies,
keeps an approximately fixed coordinate label even though its physical
distance from another comoving object changes.

\begin{figure}[H]
  \centering
  \includegraphics[width=0.72\textwidth]{lecture_09_figure_04.png}
  \caption{A coordinate mesh carried along by an expanding spatial sheet.
  Source frame at 00:22:20.}
  \label{fig:l09-comoving-grid-board}
\end{figure}

For two objects separated along the \(x\)-axis by a fixed comoving interval
\(\Delta x\), their physical distance is
\begin{equation}
  D(t)=a(t)\,\Delta x.
  \label{eq:l09-physical-distance}
\end{equation}
Differentiating while holding \(\Delta x\) fixed gives
\begin{equation}
  v(t)=\dot a(t)\,\Delta x
      =\frac{\dot a(t)}{a(t)}D(t).
\end{equation}
Defining
\begin{equation}
  H(t)\equiv\frac{\dot a(t)}{a(t)},
  \label{eq:l09-hubble-parameter}
\end{equation}
we obtain Hubble's law,
\begin{equation}
  v=H(t)D.
  \label{eq:l09-hubble-law}
\end{equation}
The word "constant" in "Hubble constant" means constant across space at a
specified cosmic time. It does not imply constancy through cosmic history.

\begin{figure}[H]
  \centering
  \includegraphics[width=0.72\textwidth]{lecture_09_figure_05.png}
  \caption{The board derivation \(D=a\Delta x\),
  \(v=(\dot a/a)D=HD\). Source frame at 00:26:40.}
  \label{fig:l09-hubble-board}
\end{figure}

\begin{classroomqa}
\audiencequestion{If a distant galaxy is seen at an earlier time, should its
inferred Hubble parameter really equal the value measured nearby today?}
\lecturerresponse{Only when the lookback time is small compared with the age
of the universe. Over modest distances, \(H(t)\) changes little enough that one
number is useful. At truly cosmological distances we see earlier epochs, and
the time dependence of \(a\) and \(H\) must be included. That is one reason we
need a dynamical equation for the scale factor.}
\end{classroomqa}

\section{Null rays in an expanding coordinate system}

\begin{classroomqa}
\audiencequestion{Does the scale factor affect the propagation of light
between distant galaxies?}
\lecturerresponse{A light ray follows a null trajectory, so \(d\tau=0\). For a
ray moving along the positive \(x\)-axis, \(dy=dz=0\), and the expanding metric
gives
\begin{align}
  0&=dt^2-a^2(t)dx^2,
  &dt&=a(t)dx,\notag\\
  &&\frac{dx}{dt}&=\frac{1}{a(t)}.
  \label{eq:l09-null-ray}
\end{align}
The coordinate speed decreases as \(a(t)\) grows, but the local physical speed
of light remains \(c=1\). A fixed coordinate interval represents an
increasing physical distance.}
\end{classroomqa}

\begin{figure}[H]
  \centering
  \includegraphics[width=0.72\textwidth]{lecture_09_figure_06.png}
  \caption{The null condition and \(dt=a(t)dx\) for radial light propagation.
  Source frame at 00:30:54.}
  \label{fig:l09-null-board}
\end{figure}

One can choose coordinates in which this same fact sounds like "light slows
down farther away," but then lengths and clocks acquire correspondingly awkward
descriptions. The expanding-space language is cleaner because it preserves
the local statement that every freely falling observer measures the same
speed of light.

\begin{classroomqa}
\audiencequestion{Is there evidence that local light propagation and atomic
physics are the same in distant galaxies?}
\lecturerresponse{Only dimensionless comparisons are operationally
meaningful. Choose an atomic size as a unit of length and an atomic or nuclear
period as a unit of time, then ask how many such periods light requires to
cross the chosen length. Spectral observations support the statement that
these dimensionless ratios are the same in distant regions. The evidence is
indirect, but it is precisely the kind of comparison by which one tests the
uniformity of local laws.}
\end{classroomqa}

When the observable universe was extremely small, light could cross a given
physical region quickly in the corresponding light-crossing unit. This does
not mean that \(c\) was different. It means that the region itself was small.

\section{A Newtonian sphere inside the universe}

We now turn from kinematics to dynamics. Place ourselves at the origin and
draw a sphere large enough to contain many galaxies and clusters. Select a
galaxy on its boundary. In a homogeneous universe the mass distribution is
spherically symmetric around us, and Newton's shell theorem gives two crucial
facts:
\begin{enumerate}
  \item all spherically distributed matter outside the chosen radius exerts
  zero net Newtonian force on the boundary galaxy;
  \item all matter inside the sphere acts, for this radial motion, as if it
  were concentrated at the center.
\end{enumerate}
We use a frame in which we are at rest, so velocities and accelerations in the
following argument are measured relative to us.

\begin{figure}[H]
  \centering
  \includegraphics[width=0.72\textwidth]{lecture_09_figure_08.png}
  \caption{The shell-theorem construction: us at the center and a test galaxy
  on the boundary of a large homogeneous sphere. Source frame at 00:41:10.}
  \label{fig:l09-shell-board}
\end{figure}

\begin{figure}[H]
  \centering
  \begin{tikzpicture}[scale=0.92,>=Latex]
    \draw[thick] (0,0) circle (2.35);
    \fill (0,0) circle (2pt) node[below left] {us};
    \fill (2.35,0) circle (2.2pt) node[right] {test galaxy \(m\)};
    \draw[->,thick,blue!65!black] (0.18,0) -- (2.15,0)
      node[midway,above] {\(R=a(t)x\)};
    \draw[->,thick,red!65!black] (2.35,0.12) -- (3.25,0.12)
      node[midway,above] {\(v=\dot a(t)x\)};
    \node[align=center,text width=7cm] at (0,-2.85)
      {outside shells cancel;\\enclosed mass acts at the center};
  \end{tikzpicture}
  \caption{Clean reconstruction of the Newtonian construction.}
  \label{fig:l09-shell-clean}
\end{figure}

If the galaxy has fixed comoving coordinate \(x\), its physical radius and
velocity are
\begin{equation}
  R(t)=x a(t),
  \qquad
  v(t)=x\dot a(t)=H(t)R(t).
  \label{eq:l09-radius-velocity}
\end{equation}
Let \(M\) be the mass enclosed by the sphere and \(m\) the mass of the test
galaxy. The Newtonian energy is
\begin{equation}
  E=\frac12 m v^2-\frac{GMm}{R}.
  \label{eq:l09-newtonian-energy}
\end{equation}

\begin{figure}[H]
  \centering
  \includegraphics[width=0.72\textwidth]{lecture_09_figure_09.png}
  \caption{The complete board sequence \(R=xa\), \(v=x\dot a\), and the
  Newtonian energy equation specialized to \(E=0\). Source frame at 00:48:40.}
  \label{fig:l09-energy-board}
\end{figure}

The energy may be positive, zero, or negative. Start at the boundary between
escape and recollapse, \(E=0\), for which
\begin{equation}
  \frac12 v^2=\frac{GM}{R}.
  \label{eq:l09-zero-energy}
\end{equation}

\begin{classroomqa}
\audiencequestion{What does zero total energy mean for the receding galaxy?}
\lecturerresponse{It is exactly the escape-velocity case. Positive energy
means motion above escape velocity; negative energy means motion below escape
velocity and, in the matter-only Newtonian picture, eventual return. We begin
at the boundary between those possibilities and return to the other signs
after deriving the flat equation.}
\end{classroomqa}

\begin{classroomqa}
\audiencequestion{Does the enclosed mass \(M\) vary as the sphere expands?}
\lecturerresponse{No. The sphere is comoving, so the same particles remain
inside it. Its physical volume grows and its density falls, but its total
enclosed mass remains fixed. This distinction between \(M\) and \(\rho(t)\) is
essential.}
\end{classroomqa}

For a homogeneous density \(\rho(t)\),
\begin{equation}
  M=\frac{4\pi}{3}\rho(t)R^3.
  \label{eq:l09-enclosed-mass}
\end{equation}
Substitution into \eqref{eq:l09-zero-energy} gives
\begin{equation}
  \frac12 v^2
  =\frac{G}{R}\left(\frac{4\pi}{3}\rho R^3\right).
  \label{eq:l09-enclosed-energy}
\end{equation}

\begin{figure}[H]
  \centering
  \includegraphics[width=0.72\textwidth]{lecture_09_figure_10.png}
  \caption{The enclosed-mass substitution in the zero-energy equation.
  Source frame at 00:54:16.060.}
  \label{fig:l09-enclosed-mass-board}
\end{figure}

Using \(v=HR\), we find
\begin{align}
  \frac12 H^2R^2
    &=\frac{4\pi G}{3}\rho R^2,\\
  H^2
    &=\frac{8\pi G}{3}\rho.
\end{align}
Thus the zero-energy, spatially flat expansion law is
\begin{equation}
  \boxed{\left(\frac{\dot a}{a}\right)^2
  =\frac{8\pi G}{3}\rho(t)}.
  \label{eq:l09-flat-friedmann}
\end{equation}
The cancellation of \(R^2\) is not accidental. A homogeneous universe must
give the same \(H(t)\) no matter which comoving galaxy was chosen as the test
object.

\begin{figure}[H]
  \centering
  \includegraphics[width=0.72\textwidth]{lecture_09_figure_11.png}
  \caption{The cancellation of the radius and the resulting flat Friedmann
  relation on the board. Source frame at 00:57:30.}
  \label{fig:l09-friedmann-board}
\end{figure}

\begin{classroomqa}
\audiencequestion{Is there evidence for the zero-energy relation?}
\lecturerresponse{Both sides can be inferred observationally: \(H\) from
recession data and \(\rho\) by accounting for stars, gas, dark matter, and
vacuum energy. The lecture-era comparison found the two sides consistent to
about percent accuracy. The qualification matters: all forms of energy must
be included, and the later vacuum-energy discussion changes the naive
escape-velocity interpretation.}
\end{classroomqa}

In units \(c=1\), mass density and energy density have the same units. The
symbol \(\rho\) will therefore denote the total energy density appropriate to
the component under discussion.

\section{Ordinary matter and the two-thirds law}

Equation \eqref{eq:l09-flat-friedmann} is not closed until we know how
\(\rho(t)\) changes. First fill a comoving rectangular box with slowly moving,
nonrelativistic particles. Its coordinate dimensions are fixed, while its
physical volume grows as
\begin{equation}
  V(t)=a^3(t)\,\Delta x\,\Delta y\,\Delta z.
\end{equation}
The particle number and rest mass in the box remain fixed, so
\begin{equation}
  \rho_{\mathrm m}(t)=\frac{C_{\mathrm m}}{a^3(t)}.
  \label{eq:l09-matter-dilution}
\end{equation}
Normalizing to the present epoch \(t_0\),
\begin{equation}
  \frac{\rho_{\mathrm m}(t)}{\rho_{\mathrm m,0}}
  =\left(\frac{a_0}{a(t)}\right)^3.
  \label{eq:l09-matter-normalized}
\end{equation}

\begin{figure}[H]
  \centering
  \includegraphics[width=0.72\textwidth]{lecture_09_figure_12.png}
  \caption{Matter density, its \(a^{-3}\) dilution, and present-day
  normalization. Source frame at 01:04:47.940.}
  \label{fig:l09-density-board}
\end{figure}

\begin{classroomqa}
\audiencequestion{What is the rough present energy density of the universe,
and can ordinary atoms alone provide it?}
\lecturerresponse{A count of visible baryonic matter gives only part of the
answer. Dark matter and dark energy must also be represented as equivalent
mass-energy. The lecture quotes a rough scale ranging from an atom to a few
dozen hydrogen-mass equivalents per cubic meter, emphasizing order of
magnitude rather than a precision inventory.}
\end{classroomqa}

Insert \eqref{eq:l09-matter-dilution} into
\eqref{eq:l09-flat-friedmann} and absorb all time-independent factors into a
positive constant \(C\):
\begin{equation}
  \left(\frac{\dot a}{a}\right)^2=\frac{C}{a^3},
  \qquad
  \dot a^2=\frac{C}{a}.
  \label{eq:l09-matter-ode}
\end{equation}
There are two transparent ways to solve it. A power-law trial
\(a=K t^p\) requires
\begin{equation}
  p-1=-\frac{p}{2},
  \qquad p=\frac23.
\end{equation}
Equivalently, take the expanding square-root branch and integrate:
\begin{equation}
  a^{1/2}\frac{da}{dt}=\sqrt C,
  \qquad
  \int a^{1/2}da=\sqrt C\int dt,
\end{equation}
so, after choosing the Big Bang origin of time,
\begin{equation}
  a^{3/2}\propto t,
  \qquad
  \boxed{a(t)\propto t^{2/3}}.
  \label{eq:l09-matter-solution}
\end{equation}
The live derivation briefly loses a factor while matching coefficients, then
returns to the integral and verifies the exponent. The notes retain the check
but display only the corrected algebra.

\begin{figure}[H]
  \centering
  \includegraphics[width=0.72\textwidth]{lecture_09_figure_13.png}
  \caption{The power-law trial and the boxed exponent \(p=2/3\). Source frame
  at 01:13:05.}
  \label{fig:l09-matter-solution-board}
\end{figure}

Differentiating the solution gives
\begin{equation}
  H(t)=\frac{\dot a}{a}=\frac{2}{3t}.
  \label{eq:l09-matter-hubble}
\end{equation}
The Hubble parameter was larger in the past. Consequently, when looking far
enough away to see appreciable lookback time, one should not extend the local
linear relation using today's \(H_0\) unchanged.

\begin{figure}[H]
  \centering
  \includegraphics[width=0.72\textwidth]{lecture_09_figure_14.png}
  \caption{Taking \(\dot a/a\) for \(a\propto t^{2/3}\) gives
  \(H=2/(3t)\). Source frame at 01:20:42.}
  \label{fig:l09-hubble-time-board}
\end{figure}

\section{Energy sign and spatial curvature}

Newtonian reasoning has carried us through the flat zero-energy case. General
relativity adds a fact that cannot be obtained from the shell theorem alone:
the sign of the Newtonian-style specific energy corresponds to the sign of
the spatial curvature.
\begin{center}
\begin{tabular}{@{}>{\raggedright\arraybackslash}p{0.20\linewidth}
                    >{\raggedright\arraybackslash}p{0.29\linewidth}
                    >{\raggedright\arraybackslash}p{0.37\linewidth}@{}}
  \toprule
  Total energy & Spatial curvature & Matter-only fate \\
  \midrule
  positive & negative (open) & expands forever \\
  zero & zero (flat) & approaches escape boundary \\
  negative & positive (closed) & expands, turns, recollapses \\
  \bottomrule
\end{tabular}
\end{center}
The last column assumes no vacuum energy. In the projectile analogy, positive
energy means above escape velocity, zero means exactly at escape velocity,
and negative means below escape velocity. Einstein's equations identify these
three cases with open, flat, and closed spatial geometry respectively.

\section{Photons and the one-half law}

The matter law \(\rho\propto a^{-3}\) assumes that the energy of each particle
is approximately fixed while the number density falls. Photons violate that
assumption. For a photon,
\begin{equation}
  E=h\nu=\frac{hc}{\lambda}.
\end{equation}
Expansion stretches its wavelength in proportion to the scale factor,
\begin{equation}
  \lambda\propto a,
  \qquad E\propto a^{-1}.
\end{equation}
The photon number density still falls as \(a^{-3}\), and each photon's energy
contributes one additional factor of \(a^{-1}\). Therefore
\begin{equation}
  \boxed{\rho_{\mathrm r}(t)\propto a^{-4}(t)}.
  \label{eq:l09-radiation-dilution}
\end{equation}

The guitar-string analogy captures the physical picture. Pluck a string and
then lengthen its vibrating portion: the wavelength lengthens and the tone
drops. A photon mode in an expanding comoving box behaves similarly. Running
the picture backward, compression shortens wavelengths and raises photon
energies.

\begin{figure}[H]
  \centering
  \includegraphics[width=0.72\textwidth]{lecture_09_figure_15.png}
  \caption{The matter-density relation extended on the board to the radiation
  scaling \(\rho_{\mathrm r}\propto a^{-4}\). Source frame at 01:29:20.}
  \label{fig:l09-radiation-board}
\end{figure}

Using \(\rho_{\mathrm r}\propto a^{-4}\) in the flat Friedmann equation gives
\begin{equation}
  \left(\frac{\dot a}{a}\right)^2=\frac{C_{\mathrm r}}{a^4},
  \qquad
  \dot a\propto a^{-1},
\end{equation}
and hence
\begin{equation}
  \boxed{a(t)\propto t^{1/2}}.
  \label{eq:l09-radiation-solution}
\end{equation}
This grows more slowly than \(t^{2/3}\). Far enough into the past, radiation
must dominate ordinary matter because \(a^{-4}\) rises faster than \(a^{-3}\)
as \(a\) decreases. The lecture places the matter-radiation crossover around
the epoch when the scale factor was roughly a thousand times smaller than
today. Before that crossover the universe was radiation dominated; after it,
for a long interval, it was matter dominated.

\section{Vacuum energy and exponential expansion}

Vacuum energy has the opposite scaling behavior. Its density is a property of
empty space and does not dilute as the universe expands:
\begin{equation}
  \rho_\Lambda=\text{constant}.
  \label{eq:l09-vacuum-density}
\end{equation}
Ordinary matter and radiation become denser toward the past, while this term
does not. It was therefore less important in the early universe and becomes
more important as the other components dilute. The lecture uses the
then-current estimate that vacuum energy supplies about \(70\%\) of today's
cosmic energy density.

With constant \(\rho_\Lambda\), equation \eqref{eq:l09-flat-friedmann} gives a
constant expansion rate
\begin{equation}
  H_\Lambda=\sqrt{\frac{8\pi G}{3}\rho_\Lambda},
  \qquad
  \frac{\dot a}{a}=H_\Lambda.
\end{equation}
Therefore
\begin{equation}
  \boxed{a(t)=a(t_i)e^{H_\Lambda(t-t_i)}}.
  \label{eq:l09-vacuum-solution}
\end{equation}

\begin{figure}[H]
  \centering
  \includegraphics[width=0.72\textwidth]{lecture_09_figure_16.png}
  \caption{The constant-density equation leading to exponential expansion.
  Source frame at 01:35:48.}
  \label{fig:l09-vacuum-board}
\end{figure}

At the observed scale quoted in the lecture, the universe doubles in linear
size over roughly ten to twelve billion years. The origin of this small but
nonzero vacuum density remains unexplained; its existence and its effect on
the expansion are empirical facts.

\begin{classroomqa}
\audiencequestion{Is vacuum energy the same thing as dark energy in this
discussion?}
\lecturerresponse{Yes. Here "dark energy" denotes a constant vacuum-energy
density, equivalently a cosmological constant. More general time-dependent
dark-energy models can be imagined, but they are not part of this lecture's
calculation.}
\end{classroomqa}

\section{How an expanding universe redshifts light}

\begin{classroomqa}
\audiencequestion{How can photons know that we chose expanding comoving
coordinates and therefore lower their energy?}
\lecturerresponse{They do not respond to our labels. They propagate through a
geometry whose physical separations change. A wave train occupying a region
that doubles in physical size acquires a doubled wavelength. The comoving
coordinates merely describe this stretching economically; they do not cause
it.}
\end{classroomqa}

\begin{classroomqa}
\audiencequestion{Where do the background photons reaching us today come
from?}
\lecturerresponse{Looking farther away means looking earlier. Eventually we
see the epoch when the cosmic material was an ionized plasma, opaque because
photons scattered repeatedly from free charged particles. The boundary from
which photons could last travel freely to us appears as a sphere surrounding
the observer: the surface of last scattering.}
\end{classroomqa}

\begin{figure}[H]
  \centering
  \includegraphics[width=0.72\textwidth]{lecture_09_figure_17.png}
  \caption{The observer-centered sketch of the surface of last scattering and
  increasingly distant shells. Source frame at 01:42:05.}
  \label{fig:l09-last-scattering-board}
\end{figure}

\begin{figure}[H]
  \centering
  \begin{tikzpicture}[scale=0.9,>=Latex]
    \fill (0,0) circle (2pt) node[below] {observer};
    \draw[thick,orange!70!black] (0,0) circle (2.55);
    \node[orange!70!black,align=center] at (0,3.05)
      {surface of last scattering};
    \foreach \a in {25,75,145,210,285,335}{
      \draw[->,blue!60!black] ({2.55*cos(\a)},{2.55*sin(\a)}) --
        ({1.7*cos(\a)},{1.7*sin(\a)});
    }
    \draw[decorate,decoration={snake,amplitude=1.5pt,segment length=7pt},
      red!65!black,thick,->] (-2.5,0) -- (-0.15,0);
    \node[red!65!black,below] at (-1.35,-0.15)
      {redshifted photon};
  \end{tikzpicture}
  \caption{Clean reconstruction of the last-scattering picture.}
  \label{fig:l09-last-scattering-clean}
\end{figure}

The radiation was emitted with a spectrum characteristic of a hot plasma and
reaches us with wavelengths roughly a thousand times longer. In the
lecture's Doppler language, the distant emitting surface recedes rapidly and
its radiation is strongly redshifted. In the geometric language, the wave is
stretched during propagation. Both pictures encode the same relation,
\begin{equation}
  1+z=\frac{\lambda_0}{\lambda_{\mathrm{em}}}
      =\frac{a_0}{a_{\mathrm{em}}}.
  \label{eq:l09-cosmological-redshift}
\end{equation}
\editorialnote{Equation \eqref{eq:l09-cosmological-redshift} is the standard
quantitative form of the lecture's stretching and Doppler descriptions.}

At an earlier observing epoch, the last-scattering sphere would have been
closer in the corresponding picture, and the observed background photons
would have been less redshifted and therefore more energetic. As cosmic time
advances, photons reaching a comoving observer from that early plasma have
progressively longer wavelengths.

\section{Where does a redshifted photon's energy go?}

\begin{classroomqa}
\audiencequestion{A photon is emitted with energy \(h\nu_{\mathrm{em}}\) and
arrives redshifted with less energy. What happened to the difference?}
\lecturerresponse{The Newtonian derivation began as an energy equation. Its
cosmological form relates the expansion rate to the material energy density,
so matter energy and expansion cannot be treated as separately conserved
accounts. The lecture describes the lost photon energy as entering the
kinetic bookkeeping of expansion. A piston filled with radiation makes the
analogy concrete: radiation pressure pushes the piston outward, the photon
wavelengths increase, and their lost energy becomes mechanical work and
kinetic energy of the piston.}
\end{classroomqa}

\begin{figure}[H]
  \centering
  \includegraphics[width=0.72\textwidth]{lecture_09_figure_18.png}
  \caption{The boxed Newtonian energy equation above its Friedmann form during
  the energy-conservation discussion. Source frame at 01:44:42.440.}
  \label{fig:l09-energy-bookkeeping-board}
\end{figure}

\begin{classroomqa}
\audiencequestion{What couples the radiation energy to the expansion?}
\lecturerresponse{Pressure. In a cylinder with a movable piston, photons
transfer momentum to the wall. As the piston moves, the radiation does work;
its wavelength grows and its energy falls. The homogeneous expanding universe
has no literal outer piston, but the \(p\,dV\) analogy explains why the energy
density of radiation falls faster than number dilution alone.}
\end{classroomqa}

\begin{figure}[H]
  \centering
  \includegraphics[width=0.72\textwidth]{lecture_09_figure_19.png}
  \caption{The movable photon-filled piston used to visualize pressure work.
  Source frame at 01:46:40.}
  \label{fig:l09-piston-board}
\end{figure}

\begin{figure}[H]
  \centering
  \begin{tikzpicture}[scale=0.92,>=Latex]
    \draw[thick] (0,0) rectangle (5.3,2.6);
    \draw[very thick] (4.25,0) -- (4.25,2.6);
    \draw[very thick] (4.25,1.3) -- (5.25,1.3);
    \draw[->,thick,red!65!black] (5.25,1.3) -- (6.2,1.3)
      node[right] {moving piston};
    \foreach \y in {0.55,1.3,2.05}{
      \draw[decorate,decoration={snake,amplitude=1.4pt,segment length=6pt},
        blue!65!black,->] (0.35,\y) -- (3.9,\y);
    }
    \node at (2.15,-0.45) {radiation does \(p\,dV\) work};
  \end{tikzpicture}
  \caption{Clean reconstruction of the radiation-piston analogy.}
  \label{fig:l09-piston-clean}
\end{figure}

Energy in general relativity requires care. In a generic expanding spacetime
there need not be a unique conserved global energy analogous to the energy of
an isolated Newtonian system. Local stress-energy conservation still holds,
and for a homogeneous fluid it yields the continuity equation
\begin{equation}
  \dot\rho+3H(\rho+p)=0,
  \label{eq:l09-continuity}
\end{equation}
which contains the pressure-work term behind the matter and radiation scaling
laws.\editorialnote{Equation \eqref{eq:l09-continuity} is the precise local
general-relativistic statement underlying the lecture's deliberately
Newtonian energy-bookkeeping analogy.}

Vacuum energy makes the puzzle especially sharp. Since its density stays
constant while volume grows, the total vacuum energy in a comoving region
increases. Its negative pressure is part of the gravitational source and
drives accelerated expansion. One should therefore not demand separate
conservation of the naive "matter energy in a comoving box" without including
the geometry and pressure work.

\section{Why the scale factor can be removed from the time term}

\begin{classroomqa}
\audiencequestion{Why does the scale factor multiply the spatial part of the
metric but not the time part?}
\lecturerresponse{That is a choice of time coordinate. Suppose one begins in a
conformal-time coordinate \(t\) with
\begin{equation}
  d\tau^2=A^2(t)\left(dt^2-dx^2\right).
\end{equation}
Define a new cosmic time \(T\) by
\begin{equation}
  dT=A(t)\,dt.
\end{equation}
Then
\begin{equation}
  d\tau^2=dT^2-A^2\bigl(t(T)\bigr)dx^2.
\end{equation}
The common factor in front of the time differential has disappeared, while
the physical scale factor remains in the spatial term.}
\end{classroomqa}

\begin{figure}[H]
  \centering
  \includegraphics[width=0.72\textwidth]{lecture_09_figure_20.png}
  \caption{The time redefinition \(dT=A(t)dt\) that converts conformal time to
  cosmic time. Source frame at 01:51:05.}
  \label{fig:l09-time-reparam-board}
\end{figure}

\section{Curvature, vacuum energy, and cosmic fate}

\begin{classroomqa}
\audiencequestion{If the universe were positively curved, should we already
see galaxies coming toward us?}
\lecturerresponse{Not necessarily. A projectile launched below escape
velocity first travels outward, reaches a turning point, and only later falls
back. We observe the universe during an expanding phase. Without vacuum
energy, a positively curved matter universe would eventually recollapse. A
sufficiently large positive vacuum energy prevents that turnaround, and the
observed value appears large enough to make the long-term expansion
exponential regardless of a small spatial curvature.}
\end{classroomqa}

The exponential future is not an energetic paradise. Stars eventually burn
out, matter and radiation densities fall exponentially, and distant regions
pass beyond causal contact. Vacuum energy remains while ordinary structure
becomes progressively colder and more dilute.

\begin{classroomqa}
\audiencequestion{Looking farther into the past, are there more photons than
ordinary particles?}
\lecturerresponse{The photon-to-baryon number ratio is already enormous,
roughly \(10^{10}\) photons per baryon in the lecture's estimate, and both
numbers in a comoving volume remain approximately fixed. What changes toward
the past is mainly the energy per photon. Photon energies rise as \(a^{-1}\),
while nonrelativistic rest masses do not. Radiation therefore dominates the
energy density at early times even though the number ratio itself need not
change.}
\end{classroomqa}

The central lesson is that geometry alone does not determine the cosmic
history. The flat Friedmann equation supplies the relation between expansion
and energy density, but the material content supplies the time dependence.
Nonrelativistic matter gives \(a\propto t^{2/3}\), radiation gives
\(a\propto t^{1/2}\), and constant vacuum energy gives exponential expansion.
Those three laws organize the successive eras of the universe and explain why
the same scale factor controls distances, redshifts, and the changing Hubble
parameter.
